\documentclass[12pt, a4paper]{article}

% ──────────────────────────────────────────────
% PAQUETES
% ──────────────────────────────────────────────
\usepackage[utf8]{inputenc}
\usepackage[T1]{fontenc}
\usepackage[spanish]{babel}
\usepackage{geometry}
\geometry{margin=2.5cm}
\usepackage{graphicx}
\usepackage{xcolor}
\usepackage{tikz}
\usetikzlibrary{shapes.geometric, arrows.meta, positioning, shadows, calc, fit}
\usepackage{enumitem}
\usepackage{fancyhdr}
\usepackage{tcolorbox}
\tcbuselibrary{skins, breakable}
\usepackage{fontawesome5}
\usepackage{hyperref}
\usepackage{tabularx}
\usepackage{booktabs}
\usepackage{multicol}
\usepackage{setspace}
\usepackage{titlesec}

% ──────────────────────────────────────────────
% COLORES YEDATECH
% ──────────────────────────────────────────────
\definecolor{ytOrange}{HTML}{FF6900}
\definecolor{ytDark}{HTML}{1A1A2E}
\definecolor{ytGrey}{HTML}{2A2A2A}
\definecolor{ytLightGrey}{HTML}{F5F5F5}
\definecolor{ytGreen}{HTML}{10B981}
\definecolor{ytRed}{HTML}{EF4444}
\definecolor{ytBlue}{HTML}{6366F1}
\definecolor{ytYellow}{HTML}{F59E0B}

% ──────────────────────────────────────────────
% ESTILOS
% ──────────────────────────────────────────────
\hypersetup{
    colorlinks=true,
    linkcolor=ytOrange,
    urlcolor=ytBlue,
}

\pagestyle{fancy}
\fancyhf{}
\fancyhead[L]{\textcolor{ytOrange}{\textbf{YedaTech}} \textcolor{ytGrey}{— Manual del Usuario}}
\fancyhead[R]{\textcolor{ytGrey}{\thepage}}
\fancyfoot[C]{\small\textcolor{gray}{Tutor Inteligente de Matemáticas con IA}}
\renewcommand{\headrulewidth}{0.5pt}
\renewcommand{\headrule}{\hbox to\headwidth{\color{ytOrange}\leaders\hrule height \headrulewidth\hfill}}

\titleformat{\section}
  {\Large\bfseries\color{ytDark}}{\textcolor{ytOrange}{\thesection}}{1em}{}
  [\vspace{-0.5em}\textcolor{ytOrange}{\rule{\textwidth}{1pt}}]

\titleformat{\subsection}
  {\large\bfseries\color{ytGrey}}{}{0em}{\textcolor{ytOrange}{\faIcon{chevron-right}}\ }

% ──────────────────────────────────────────────
% CAJAS PERSONALIZADAS
% ──────────────────────────────────────────────
\newtcolorbox{infobox}[1][]{
    enhanced, breakable,
    colback=ytOrange!5, colframe=ytOrange,
    fonttitle=\bfseries\color{white},
    title={#1},
    coltitle=white,
    attach boxed title to top left={yshift=-2mm, xshift=5mm},
    boxed title style={colback=ytOrange, rounded corners},
    rounded corners,
    drop shadow,
}

\newtcolorbox{tipbox}[1][]{
    enhanced, breakable,
    colback=ytGreen!5, colframe=ytGreen,
    fonttitle=\bfseries\color{white},
    title={#1},
    coltitle=white,
    attach boxed title to top left={yshift=-2mm, xshift=5mm},
    boxed title style={colback=ytGreen, rounded corners},
    rounded corners,
}

\newtcolorbox{warningbox}[1][]{
    enhanced, breakable,
    colback=ytYellow!5, colframe=ytYellow,
    fonttitle=\bfseries\color{white},
    title={#1},
    coltitle=white,
    attach boxed title to top left={yshift=-2mm, xshift=5mm},
    boxed title style={colback=ytYellow, rounded corners},
    rounded corners,
}

% ──────────────────────────────────────────────
% DOCUMENTO
% ──────────────────────────────────────────────
\begin{document}

% ══════════════════════════════════════════════
% PORTADA
% ══════════════════════════════════════════════
\begin{titlepage}
\begin{center}

\vspace*{2cm}

% Logo simulado
\begin{tikzpicture}
    \node[fill=ytOrange, rounded corners=12pt, minimum width=2.5cm, minimum height=2.5cm, text=white, font=\Huge\bfseries] {YT};
\end{tikzpicture}

\vspace{1cm}

{\Huge\bfseries\textcolor{ytDark}{YedaTech}}\\[0.3cm]
{\LARGE\textcolor{ytOrange}{Tutor Inteligente de Matemáticas}}\\[0.2cm]
{\large\textcolor{ytGrey}{Potenciado por Inteligencia Artificial}}

\vspace{1.5cm}

\textcolor{ytOrange}{\rule{0.6\textwidth}{2pt}}

\vspace{1.5cm}

{\Large\textcolor{ytGrey}{\textbf{Manual del Usuario}}}\\[0.5cm]
{\large\textcolor{gray}{Guía completa del sistema para estudiantes y docentes}}

\vspace{2cm}

\begin{tikzpicture}
    \node[draw=ytOrange!30, rounded corners=8pt, fill=ytOrange!3, inner sep=15pt, text width=10cm, align=center] {
        \textcolor{ytGrey}{
            \faIcon{graduation-cap}\ \ Evaluación diagnóstica personalizada\\[3pt]
            \faIcon{brain}\ \ Tutoría socrática adaptativa\\[3pt]
            \faIcon{chart-line}\ \ Analítica de progreso en tiempo real\\[3pt]
            \faIcon{book}\ \ Alineado al currículo del MEN de Colombia
        }
    };
\end{tikzpicture}

\vfill

{\small\textcolor{gray}{Versión 1.0 — Marzo 2026}}

\end{center}
\end{titlepage}

% ══════════════════════════════════════════════
% TABLA DE CONTENIDOS
% ══════════════════════════════════════════════
\tableofcontents
\newpage

% ══════════════════════════════════════════════
% 1. ¿QUÉ ES YEDATECH?
% ══════════════════════════════════════════════
\section{¿Qué es YedaTech?}

YedaTech es una plataforma educativa inteligente diseñada para ayudar a estudiantes de bachillerato colombiano (grados 6° a 11°) a mejorar sus habilidades en matemáticas. A diferencia de una clase tradicional, YedaTech utiliza \textbf{Inteligencia Artificial} para adaptarse al nivel y ritmo de cada estudiante.

\begin{infobox}[\faIcon{lightbulb}\ Filosofía del Sistema]
YedaTech se basa en el \textbf{Método Socrático}: en lugar de dar respuestas directas, el tutor guía al estudiante con preguntas estratégicas para que descubra la solución por sí mismo. Esto promueve el pensamiento crítico y la comprensión profunda.
\end{infobox}

\subsection{Características principales}

\begin{itemize}[leftmargin=*, itemsep=6pt]
    \item[\textcolor{ytOrange}{\faIcon{clipboard-check}}] \textbf{Evaluación Diagnóstica Inteligente:} Un examen de 15 preguntas generado por IA que identifica fortalezas y debilidades del estudiante.
    \item[\textcolor{ytOrange}{\faIcon{brain}}] \textbf{Tutor Socrático con IA:} Un chat interactivo donde el tutor guía al estudiante paso a paso sin revelar la respuesta.
    \item[\textcolor{ytOrange}{\faIcon{sliders-h}}] \textbf{Dificultad Adaptativa:} El sistema ajusta automáticamente la complejidad de los ejercicios según el desempeño del estudiante.
    \item[\textcolor{ytOrange}{\faIcon{chart-bar}}] \textbf{Analítica Detallada:} Gráficas y reportes que muestran el progreso por tema, tipo de error, y rendimiento temporal.
    \item[\textcolor{ytOrange}{\faIcon{flag}}] \textbf{Currículo Oficial MEN:} Todos los contenidos están alineados con los Estándares Básicos de Competencias y los DBA del Ministerio de Educación Nacional de Colombia.
\end{itemize}

\subsection{¿A quién está dirigido?}

\begin{center}
\begin{tabularx}{0.9\textwidth}{>{\centering\arraybackslash}X >{\centering\arraybackslash}X >{\centering\arraybackslash}X}
\toprule
\textbf{\textcolor{ytOrange}{\faIcon{user-graduate}} Estudiantes} & \textbf{\textcolor{ytBlue}{\faIcon{chalkboard-teacher}} Docentes} & \textbf{\textcolor{ytGreen}{\faIcon{users}} Padres} \\
\midrule
Grados 6° a 11° de bachillerato colombiano & Para monitorear el progreso de sus estudiantes & Para acompañar el aprendizaje de sus hijos \\
\bottomrule
\end{tabularx}
\end{center}

% ══════════════════════════════════════════════
% 2. PRIMEROS PASOS
% ══════════════════════════════════════════════
\section{Primeros Pasos}

\subsection{Crear una cuenta}

Para comenzar a usar YedaTech, el estudiante debe registrarse:

\begin{enumerate}[leftmargin=*, itemsep=4pt]
    \item Acceder a la plataforma desde el navegador web.
    \item Hacer clic en \textbf{``Regístrate aquí''} en la pantalla de inicio.
    \item Completar el formulario con:
    \begin{itemize}
        \item \textbf{Nombre completo} del estudiante
        \item \textbf{Nombre de usuario} (para iniciar sesión)
        \item \textbf{Correo electrónico}
        \item \textbf{Contraseña} (mínimo 6 caracteres)
    \end{itemize}
    \item Hacer clic en \textbf{``Crear Cuenta''}.
    \item El sistema redirigirá automáticamente a la pantalla de inicio de sesión.
\end{enumerate}

\subsection{Iniciar sesión}

Una vez registrado, el estudiante inicia sesión con su nombre de usuario y contraseña. Al ingresar, verá el \textbf{Dashboard} (tablero principal) con tres opciones principales:

\begin{center}
\begin{tikzpicture}[
    box/.style={draw=ytOrange!40, rounded corners=8pt, fill=ytOrange!3, 
                minimum width=4cm, minimum height=2.2cm, align=center,
                font=\small, text width=3.5cm},
    >=Stealth
]
    \node[box] (diag) at (0,0) {\textcolor{ytOrange}{\faIcon{clipboard-check}}\\[3pt]\textbf{Evaluación}\\\textbf{Diagnóstica}\\[2pt]\scriptsize Conoce tu nivel};
    \node[box] (tutor) at (5,0) {\textcolor{ytOrange}{\faIcon{brain}}\\[3pt]\textbf{Tutor}\\\textbf{Socrático}\\[2pt]\scriptsize Practica con IA};
    \node[box] (anal) at (10,0) {\textcolor{ytOrange}{\faIcon{chart-line}}\\[3pt]\textbf{Analítica}\\\textbf{Avanzada}\\[2pt]\scriptsize Ve tu progreso};
\end{tikzpicture}
\end{center}

\begin{tipbox}[\faIcon{star}\ Recomendación]
Se recomienda \textbf{comenzar siempre por la Evaluación Diagnóstica}. Esta prueba permite al sistema conocer tu nivel actual y personalizar toda tu experiencia de aprendizaje.
\end{tipbox}

% ══════════════════════════════════════════════
% 3. EVALUACIÓN DIAGNÓSTICA
% ══════════════════════════════════════════════
\section{Evaluación Diagnóstica Inteligente}

La evaluación diagnóstica es el corazón del sistema. Es un examen personalizado que identifica exactamente qué sabe y qué necesita reforzar cada estudiante.

\subsection{¿Cómo funciona?}

El proceso se divide en tres fases:

\begin{center}
\begin{tikzpicture}[
    phase/.style={draw=ytOrange, rounded corners=10pt, fill=white, 
                  minimum width=3.8cm, minimum height=2.8cm, align=center,
                  font=\small, text width=3.5cm, drop shadow},
    arrow/.style={-{Stealth[length=8pt]}, thick, ytOrange},
]
    \node[phase] (p1) at (0,0) {
        \textcolor{ytOrange}{\Large\faIcon{cog}}\\[5pt]
        \textbf{Fase 1}\\
        \textbf{Generación}\\[3pt]
        \scriptsize La IA crea 15 preguntas únicas adaptadas a tu grado
    };
    \node[phase] (p2) at (5.5,0) {
        \textcolor{ytBlue}{\Large\faIcon{pen}}\\[5pt]
        \textbf{Fase 2}\\
        \textbf{Examen}\\[3pt]
        \scriptsize Respondes las preguntas a tu ritmo con temporizador
    };
    \node[phase] (p3) at (11,0) {
        \textcolor{ytGreen}{\Large\faIcon{chart-pie}}\\[5pt]
        \textbf{Fase 3}\\
        \textbf{Análisis}\\[3pt]
        \scriptsize 3 agentes de IA analizan tus resultados
    };
    
    \draw[arrow] (p1) -- (p2);
    \draw[arrow] (p2) -- (p3);
\end{tikzpicture}
\end{center}

\subsection{Selección de grado}

Antes de generar el examen, el estudiante selecciona su grado escolar actual (6° a 11°). Esto es fundamental porque \textbf{el examen se genera específicamente con los temas del currículo oficial} de ese grado.

\subsection{El examen}

El examen consta de \textbf{15 preguntas de opción múltiple} (A, B, C, D) distribuidas así:

\begin{center}
\begin{tabularx}{0.85\textwidth}{Xcc}
\toprule
\textbf{Pensamiento Matemático} & \textbf{Preguntas} & \textbf{Ícono} \\
\midrule
Pensamiento Numérico (Sistemas Numéricos) & 4 & \textcolor{ytOrange}{\faIcon{sort-numeric-up}} \\
Pensamiento Variacional (Álgebra y Análisis) & 4 & \textcolor{ytBlue}{\faIcon{superscript}} \\
Pensamiento Espacial y Métrico (Geometría) & 4 & \textcolor{ytGreen}{\faIcon{shapes}} \\
Pensamiento Aleatorio (Datos y Probabilidad) & 3 & \textcolor{ytYellow}{\faIcon{dice}} \\
\midrule
\textbf{Total} & \textbf{15} & \\
\bottomrule
\end{tabularx}
\end{center}

\begin{warningbox}[\faIcon{clock}\ Sobre el temporizador]
El sistema registra cuánto tiempo tardas en cada pregunta. Esto \textbf{no es para presionarte}, sino para que la IA pueda detectar si necesitas más práctica en algún tema (por ejemplo, si tardas mucho puede indicar inseguridad en ese tema).
\end{warningbox}

Durante el examen, el estudiante puede:
\begin{itemize}[itemsep=3pt]
    \item Navegar libremente entre preguntas (adelante y atrás).
    \item Ver un mapa visual de qué preguntas ya respondió.
    \item Cambiar sus respuestas antes de enviar.
    \item Enviar el examen solo cuando haya respondido \textbf{todas} las preguntas.
\end{itemize}

\subsection{Análisis de resultados}

Una vez enviado el examen, \textbf{tres agentes de inteligencia artificial} trabajan en secuencia para analizar las respuestas:

\begin{enumerate}[leftmargin=*, itemsep=8pt]
    \item \textbf{\textcolor{ytOrange}{Agente Analizador de Respuestas:}} Revisa cada pregunta incorrecta e identifica:
    \begin{itemize}
        \item El \textit{tipo de error} cometido (conceptual, procedimental, aritmético o de interpretación)
        \item El \textit{concepto débil} asociado
        \item Una \textit{explicación} de por qué la respuesta es incorrecta y cuál es la correcta
    \end{itemize}
    
    \item \textbf{\textcolor{ytBlue}{Agente de Diagnóstico Académico:}} Genera un diagnóstico global que incluye:
    \begin{itemize}
        \item Temas dominados vs.\ temas con dificultad
        \item Rendimiento por cada pensamiento matemático (Numérico, Variacional, Espacial, Aleatorio)
        \item Brecha académica respecto a lo esperado para el grado
        \item Observaciones sobre el comportamiento temporal
    \end{itemize}
    
    \item \textbf{\textcolor{ytGreen}{Agente de Recomendaciones:}} Crea un plan de estudio personalizado con:
    \begin{itemize}
        \item Temas organizados por prioridad (alta, media, baja)
        \item Ejercicios específicos recomendados para cada tema
        \item Duración estimada del plan
        \item Consejos generales de estudio
    \end{itemize}
\end{enumerate}

\subsection{Repetir el diagnóstico}

El estudiante puede \textbf{repetir la evaluación diagnóstica} en cualquier momento para medir su progreso. Al repetirla:
\begin{itemize}
    \item Se generan \textbf{preguntas nuevas} (nunca se repite el mismo examen).
    \item Se mantienen los \textbf{mismos temas} del grado para que la comparación sea válida.
    \item Se puede comparar el rendimiento anterior con el actual.
\end{itemize}

% ══════════════════════════════════════════════
% 4. TUTOR SOCRÁTICO
% ══════════════════════════════════════════════
\section{Tutor Socrático con IA}

El Tutor Socrático es la herramienta de práctica diaria. Funciona como un chat inteligente donde el estudiante resuelve ejercicios con la guía del tutor.

\subsection{¿Qué es el Método Socrático?}

El método socrático consiste en \textbf{enseñar a través de preguntas}, no de respuestas. En lugar de decirle al estudiante ``la respuesta es 42'', el tutor le hace preguntas como:

\begin{infobox}[\faIcon{comments}\ Ejemplo de diálogo socrático]
\textbf{Tutor:} \textit{Resuelve: $2x + 6 = 14$}\\[3pt]
\textbf{Estudiante:} \textit{No sé por dónde empezar}\\[3pt]
\textbf{Tutor:} \textit{¿Qué operación podrías hacer para dejar la $x$ sola en un lado? Piensa en qué está ``acompañando'' a la $x$ en este momento...}\\[3pt]
\textbf{Estudiante:} \textit{¿Restar 6 de ambos lados?}\\[3pt]
\textbf{Tutor:} \textit{¡Exacto! Y si restas 6 de ambos lados, ¿qué te queda?}
\end{infobox}

\subsection{Flujo de una sesión de tutoría}

\begin{center}
\begin{tikzpicture}[
    node distance=1.2cm,
    step/.style={draw=ytOrange!50, rounded corners=6pt, fill=ytOrange!5, 
                 minimum width=10cm, minimum height=1cm, align=center, font=\small},
    arrow/.style={-{Stealth[length=6pt]}, thick, ytOrange!60},
]
    \node[step] (s1) {\faIcon{play}\ El sistema presenta un \textbf{ejercicio} adaptado al nivel del estudiante};
    \node[step, below=of s1] (s2) {\faIcon{pen}\ El estudiante escribe su \textbf{respuesta o razonamiento}};
    \node[step, below=of s2] (s3) {\faIcon{robot}\ La IA \textbf{evalúa} si la respuesta es correcta o incorrecta};
    \node[step, below=of s3] (s4) {\faIcon{comments}\ El tutor responde con \textbf{preguntas guía} (no da la respuesta directa)};
    \node[step, below=of s4] (s5) {\faIcon{redo}\ Se repite hasta que el estudiante \textbf{descubra la respuesta} por sí mismo};
    \node[step, below=of s5] (s6) {\faIcon{arrow-up}\ Se pasa al \textbf{siguiente ejercicio} y el nivel se ajusta automáticamente};
    
    \draw[arrow] (s1) -- (s2);
    \draw[arrow] (s2) -- (s3);
    \draw[arrow] (s3) -- (s4);
    \draw[arrow] (s4) -- (s5);
    \draw[arrow] (s5) -- (s6);
\end{tikzpicture}
\end{center}

\subsection{Sistema de niveles adaptativo}

El sistema monitorea las últimas 5 respuestas evaluadas del estudiante y ajusta la dificultad automáticamente:

\begin{center}
\begin{tabularx}{0.85\textwidth}{cXc}
\toprule
\textbf{Precisión} & \textbf{Acción} & \textbf{Resultado} \\
\midrule
$\geq 80\%$ & Se \textcolor{ytGreen}{aumenta} la dificultad & \textcolor{ytGreen}{\faIcon{arrow-up}} Nivel +1 \\
$50\% - 79\%$ & Se \textcolor{ytBlue}{mantiene} la dificultad & \textcolor{ytBlue}{\faIcon{equals}} Mismo nivel \\
$< 50\%$ & Se \textcolor{ytRed}{reduce} la dificultad & \textcolor{ytRed}{\faIcon{arrow-down}} Nivel -1 \\
\bottomrule
\end{tabularx}
\end{center}

Esto garantiza que el estudiante siempre trabaje con ejercicios \textbf{desafiantes pero alcanzables}, evitando frustración o aburrimiento.

% ══════════════════════════════════════════════
% 5. ANALÍTICA Y PROGRESO
% ══════════════════════════════════════════════
\section{Analítica y Seguimiento de Progreso}

YedaTech ofrece un sistema completo de analíticas para que tanto el estudiante como el docente puedan visualizar el progreso.

\subsection{Dashboard del estudiante}

Al iniciar sesión, el estudiante ve un panel con métricas clave:

\begin{center}
\begin{tikzpicture}[
    metric/.style={draw=ytOrange!30, rounded corners=6pt, fill=white, 
                   minimum width=3.2cm, minimum height=1.8cm, align=center, font=\small, drop shadow},
]
    \node[metric] (m1) at (0,0) {\textcolor{ytOrange}{\Large\textbf{3}}\\Nivel Actual};
    \node[metric] (m2) at (4,0) {\textcolor{ytGreen}{\Large\textbf{28}}\\Respuestas\\Correctas};
    \node[metric] (m3) at (8,0) {\textcolor{ytBlue}{\Large\textbf{5}}\\Sesiones\\Completadas};
    \node[metric] (m4) at (12,0) {\textcolor{ytYellow}{\Large\textbf{73\%}}\\Precisión\\General};
\end{tikzpicture}
\end{center}

\subsection{Analítica avanzada}

La sección de analítica avanzada permite ver:

\begin{itemize}[leftmargin=*, itemsep=5pt]
    \item[\textcolor{ytOrange}{\faIcon{chart-bar}}] \textbf{Rendimiento por tema:} Porcentaje de aciertos en cada pensamiento matemático (Numérico, Variacional, Espacial/Métrico, Aleatorio).
    \item[\textcolor{ytBlue}{\faIcon{history}}] \textbf{Historial de diagnósticos:} Comparación de resultados entre evaluaciones diagnósticas para medir el avance real.
    \item[\textcolor{ytGreen}{\faIcon{clock}}] \textbf{Patrón temporal:} Cuánto tiempo promedio tarda en cada pregunta y si la velocidad mejora con el tiempo.
    \item[\textcolor{ytRed}{\faIcon{exclamation-triangle}}] \textbf{Tipos de error frecuentes:} Distribución de errores conceptuales, procedimentales, aritméticos y de interpretación.
    \item[\textcolor{ytYellow}{\faIcon{fire}}] \textbf{Rachas:} Cantidad de respuestas correctas o incorrectas consecutivas.
\end{itemize}

% ══════════════════════════════════════════════
% 6. PANEL DE ADMINISTRACIÓN
% ══════════════════════════════════════════════
\section{Panel de Administración (Docentes)}

YedaTech incluye un panel de administración exclusivo para docentes y coordinadores académicos.

\subsection{Acceso}

El panel de administración se accede con una contraseña de administrador independiente. Esto permite ver los datos de todos los estudiantes registrados en la plataforma.

\subsection{Funcionalidades del panel}

\begin{itemize}[leftmargin=*, itemsep=5pt]
    \item[\textcolor{ytOrange}{\faIcon{users}}] \textbf{Lista de Estudiantes:} Ver todos los estudiantes registrados con su nivel actual, grado, número de sesiones y última actividad.
    \item[\textcolor{ytOrange}{\faIcon{user}}] \textbf{Perfil Individual:} Al hacer clic en un estudiante, el docente puede ver:
    \begin{itemize}
        \item Su diagnóstico completo con temas dominados y débiles
        \item Historial de todas las sesiones de tutoría
        \item Detalle de cada interacción con el tutor
        \item Tipos de errores más frecuentes
    \end{itemize}
    \item[\textcolor{ytOrange}{\faIcon{chart-line}}] \textbf{Analítica Global:} Estadísticas agregadas de toda la clase, incluyendo temas más difíciles y patrones de error comunes.
\end{itemize}

% ══════════════════════════════════════════════
% 7. ARQUITECTURA DEL SISTEMA (Sin código)
% ══════════════════════════════════════════════
\section{¿Cómo funciona por dentro?}

Sin entrar en detalles técnicos de programación, YedaTech está compuesto por un equipo de \textbf{agentes inteligentes} que trabajan juntos, cada uno con una función específica:

\begin{center}
\begin{tikzpicture}[
    agent/.style={draw=#1!60, rounded corners=8pt, fill=#1!8, 
                  minimum width=4.5cm, minimum height=1.5cm, align=center,
                  font=\small, text width=4.2cm, drop shadow={shadow xshift=1pt, shadow yshift=-1pt}},
    label/.style={font=\scriptsize\bfseries, text=#1},
    conn/.style={-{Stealth[length=5pt]}, thick, #1!50},
]
    % Agentes del diagnóstico
    \node[agent=ytOrange] (gen) at (0,6) {\textcolor{ytOrange}{\faIcon{cog}}\\[2pt]\textbf{Generador de Exámenes}\\[1pt]\scriptsize Crea preguntas únicas por grado};
    \node[agent=ytRed] (ana) at (6,6) {\textcolor{ytRed}{\faIcon{search}}\\[2pt]\textbf{Analizador de Respuestas}\\[1pt]\scriptsize Identifica errores y patrones};
    \node[agent=ytBlue] (diag) at (12,6) {\textcolor{ytBlue}{\faIcon{stethoscope}}\\[2pt]\textbf{Diagnosticador Académico}\\[1pt]\scriptsize Evalúa el nivel real};
    
    % Agentes del tutor
    \node[agent=ytGreen] (router) at (0,3) {\textcolor{ytGreen}{\faIcon{route}}\\[2pt]\textbf{Clasificador de Intención}\\[1pt]\scriptsize Entiende qué dice el estudiante};
    \node[agent=ytOrange] (tutor) at (6,3) {\textcolor{ytOrange}{\faIcon{brain}}\\[2pt]\textbf{Tutor Socrático}\\[1pt]\scriptsize Guía con preguntas};
    \node[agent=ytBlue] (eval) at (12,3) {\textcolor{ytBlue}{\faIcon{check-circle}}\\[2pt]\textbf{Evaluador}\\[1pt]\scriptsize Verifica si la respuesta es correcta};
    
    % Agentes de soporte
    \node[agent=ytYellow] (adapt) at (0,0) {\textcolor{ytYellow}{\faIcon{sliders-h}}\\[2pt]\textbf{Agente Adaptativo}\\[1pt]\scriptsize Ajusta la dificultad};
    \node[agent=ytGreen] (mem) at (6,0) {\textcolor{ytGreen}{\faIcon{memory}}\\[2pt]\textbf{Agente de Memoria}\\[1pt]\scriptsize Recuerda tu perfil cognitivo};
    \node[agent=ytRed] (tel) at (12,0) {\textcolor{ytRed}{\faIcon{database}}\\[2pt]\textbf{Telemetría}\\[1pt]\scriptsize Registra toda la actividad};
    
    % Conexiones
    \draw[conn=ytOrange] (gen) -- (ana);
    \draw[conn=ytRed] (ana) -- (diag);
    \draw[conn=ytGreen] (router) -- (tutor);
    \draw[conn=ytOrange] (tutor) -- (eval);
    \draw[conn=ytBlue] (eval) -- (adapt);
    \draw[conn=ytYellow] (adapt) -- (mem);
    \draw[conn=ytGreen] (mem) -- (tel);
    
    % Labels de grupo
    \node[label=ytOrange, anchor=west] at (-2.5,6) {\rotatebox{90}{DIAGNÓSTICO}};
    \node[label=ytGreen, anchor=west] at (-2.5,3) {\rotatebox{90}{TUTORÍA}};
    \node[label=ytBlue, anchor=west] at (-2.5,0) {\rotatebox{90}{SOPORTE}};
\end{tikzpicture}
\end{center}

\begin{tipbox}[\faIcon{info-circle}\ Nota importante]
Todos estos agentes trabajan \textbf{de forma invisible} para el estudiante. El usuario solo ve las preguntas, respuestas y resultados. La complejidad del sistema está completamente oculta detrás de una interfaz sencilla e intuitiva.
\end{tipbox}

% ══════════════════════════════════════════════
% 8. CURRÍCULO
% ══════════════════════════════════════════════
\section{Alineación con el Currículo Nacional}

Todos los contenidos de YedaTech están basados en documentos oficiales del \textbf{Ministerio de Educación Nacional de Colombia (MEN)}:

\begin{itemize}[itemsep=5pt]
    \item \textbf{Lineamientos Curriculares de Matemáticas} (1998)
    \item \textbf{Estándares Básicos de Competencias en Matemáticas} (2006)
    \item \textbf{Derechos Básicos de Aprendizaje (DBA)} Versión 2 (2016)
\end{itemize}

Los contenidos se organizan en los \textbf{4 pensamientos matemáticos} definidos por el MEN:

\begin{center}
\begin{tabularx}{0.95\textwidth}{>{\bfseries}l X}
\toprule
\textbf{Pensamiento} & \textbf{¿Qué incluye?} \\
\midrule
\textcolor{ytOrange}{Numérico} & Operaciones con naturales, enteros, racionales y reales. Teoría de números. Notación científica. Matemáticas financieras (grados superiores). \\[5pt]
\textcolor{ytBlue}{Variacional} & Patrones y regularidades. Expresiones algebraicas. Ecuaciones y sistemas de ecuaciones. Funciones. Cálculo diferencial e integral (grado 11°). \\[5pt]
\textcolor{ytGreen}{Espacial y Métrico} & Geometría de polígonos y sólidos. Teorema de Pitágoras. Trigonometría. Geometría analítica. Secciones cónicas (grados superiores). \\[5pt]
\textcolor{ytYellow}{Aleatorio} & Estadística descriptiva. Medidas de tendencia central. Probabilidad. Distribuciones. Inferencia estadística (grado 11°). \\
\bottomrule
\end{tabularx}
\end{center}

\subsection{Ejemplo de contenidos por grado}

\begin{center}
\renewcommand{\arraystretch}{1.3}
\begin{tabularx}{0.95\textwidth}{c X}
\toprule
\textbf{Grado} & \textbf{Temas principales} \\
\midrule
\textbf{6°} & Operaciones con fracciones, MCD y mcm, polígonos, tablas de frecuencia \\
\textbf{7°} & Números enteros, proporcionalidad inversa, transformaciones en el plano, mediana \\
\textbf{8°} & Potenciación, factorización, productos notables, Teorema de Pitágoras \\
\textbf{9°} & Función lineal y cuadrática, sistemas de ecuaciones, semejanza de triángulos \\
\textbf{10°} & Trigonometría, secciones cónicas, interés compuesto, correlación \\
\textbf{11°} & Límites, derivadas, integrales, distribución normal, inferencia estadística \\
\bottomrule
\end{tabularx}
\end{center}

% ══════════════════════════════════════════════
% 9. PREGUNTAS FRECUENTES
% ══════════════════════════════════════════════
\section{Preguntas Frecuentes}

\begin{enumerate}[leftmargin=*, itemsep=10pt, font=\bfseries\color{ytOrange}]
    \item \textbf{¿Necesito instalar algo?}\\
    \textcolor{ytGrey}{No. YedaTech funciona directamente desde el navegador web (Chrome, Firefox, Edge, Safari). No requiere descargar ni instalar ninguna aplicación.}
    
    \item \textbf{¿El diagnóstico siempre tiene las mismas preguntas?}\\
    \textcolor{ytGrey}{No. Cada vez que se genera un examen, la inteligencia artificial crea preguntas completamente nuevas basadas en tu grado. Nunca verás el mismo examen dos veces.}
    
    \item \textbf{¿El tutor me va a dar las respuestas directamente?}\\
    \textcolor{ytGrey}{No. El tutor usa el método socrático: te guía con preguntas para que tú mismo descubras la respuesta. Si después de varios intentos no lo logras, te dará pistas cada vez más directas.}
    
    \item \textbf{¿Cómo sé si estoy mejorando?}\\
    \textcolor{ytGrey}{Puedes repetir la evaluación diagnóstica en cualquier momento para comparar resultados. Además, la sección de Analítica muestra tu progreso en gráficas y estadísticas.}
    
    \item \textbf{¿Un docente puede ver mi progreso?}\\
    \textcolor{ytGrey}{Sí, si el docente tiene acceso al panel de administración, puede ver el progreso detallado de cada estudiante, incluyendo diagnósticos, sesiones y patrones de error.}
    
    \item \textbf{¿Funciona para todos los grados?}\\
    \textcolor{ytGrey}{YedaTech cubre los grados de 6° a 11° de bachillerato colombiano, con contenidos alineados al currículo oficial del MEN para cada grado.}
    
    \item \textbf{¿Puedo usar YedaTech en el celular?}\\
    \textcolor{ytGrey}{Sí. La plataforma es responsiva y se adapta a cualquier tamaño de pantalla: computador, tablet o celular.}
\end{enumerate}

\vspace{1cm}

\begin{center}
\begin{tikzpicture}
    \node[fill=ytOrange!5, draw=ytOrange!30, rounded corners=12pt, inner sep=20pt, text width=12cm, align=center] {
        {\Large\textcolor{ytOrange}{\faIcon{graduation-cap}}}\\[8pt]
        {\large\textbf{\textcolor{ytDark}{¡Comienza tu viaje de aprendizaje con YedaTech!}}}\\[5pt]
        {\textcolor{ytGrey}{Cada ejercicio resuelto te acerca más a dominar las matemáticas.}}\\[3pt]
        {\textcolor{ytGrey}{Recuerda: no se trata de velocidad, sino de \textbf{comprensión}.}}
    };
\end{tikzpicture}
\end{center}

\end{document}
